% =====================================================================
%  Informe final: Robot de ajedrez con aprendizaje por refuerzo
%  Machine Learning, Ingeniería en IA y Ciencia de Datos, UCU (2026)
%  Compilar en Overleaf con pdfLaTeX. La carpeta figuras/ tiene que tener:
%  desde_cero_auc.jpeg, resumen_atribucion.png (la nueva, con 210k),
%  valores_piezas.png, mapa_lichess_1.png, valor_calibracion.jpeg y,
%  si la suben, una foto del montaje: montaje.jpg (es opcional: si no
%  está, el informe compila igual sin esa figura).
% =====================================================================
\documentclass[11pt,a4paper]{article}

\usepackage[utf8]{inputenc}
\usepackage[T1]{fontenc}
\usepackage[spanish,es-noshorthands,es-tabla]{babel}
\usepackage{lmodern}
\usepackage[a4paper,margin=2.3cm]{geometry}
\usepackage{amsmath,amssymb}
\usepackage{icomma}            % permite escribir 0,5 en modo matemático
\usepackage{graphicx}
\usepackage{booktabs}
\usepackage{tabularx}
\usepackage{multirow}
\usepackage{xcolor}
\usepackage{caption}
\usepackage{subcaption}
\usepackage{enumitem}
\usepackage[round,authoryear]{natbib}
\usepackage[hidelinks]{hyperref}
\usepackage[nameinlink,spanish]{cleveref}

\captionsetup{font=small,labelfont=bf}
\setlist{itemsep=1pt,topsep=3pt}
\setlength{\parskip}{3pt}
\graphicspath{{figuras/}}

\newcommand{\vhat}{\hat v_\theta}

% ---------------------------------------------------------------------
\begin{document}

% =========================== PORTADA ==================================
\begin{titlepage}
\centering
{\large Universidad Católica del Uruguay\par}
{\large Facultad de Ingeniería y Tecnologías\par}
\vspace{0.4cm}
{Ingeniería en Inteligencia Artificial y Ciencia de Datos\par}
{Machine Learning: trabajo final\par}
\vspace{2.5cm}
{\LARGE\bfseries Un robot que aprende a jugar al ajedrez\par}
\vspace{0.3cm}
{\Large Aprendizaje por refuerzo con TD-Leaf($\lambda$) y TreeStrap\par}
\vspace{2.5cm}
{\large\textbf{Integrantes}\par}
\vspace{0.2cm}
Sebastián Bascou Mortola\par
Mónica Aylín Deus Álvarez\par
Valentino Repetto Perez\par
Braihans Matihas Sastre Moreira\par
\vspace{1.5cm}
{\large\textbf{Docente:} \textit{Gustavo Vasquez y Juan Sebastián Parolli}\par}
\vfill
\begin{minipage}{0.9\textwidth}
\small
\textbf{Resumen.} Construimos un agente que aprende a jugar al ajedrez
jugando, con aprendizaje por refuerzo (RL) en una laptop, y lo pusimos
dentro de un robot que mira un tablero real con una cámara y muestra en una
pantalla sus jugadas y una cara que expresa cuán confiado está. El agente
combina una red neuronal que evalúa posiciones con una búsqueda de pocas
jugadas. Para demostrar que la mejora viene del aprendizaje comparamos
redes con todo lo demás fijo. Desde una red que solo conocía el valor de
las piezas, el RL suma +36 puntos de Elo; al estancarse, diagnosticamos la
causa, la corregimos y la red llegó a +148 Elo antes de volver a
estancarse. Un análisis de explicabilidad muestra qué aprendió: por
ejemplo, que un peón a un paso de coronar vale mucho más que uno en su
casilla inicial, algo que nadie le enseñó.
\end{minipage}
\vspace{1cm}

{\large 26 de noviembre de 2026\par}
\end{titlepage}

\setcounter{page}{1}

% =========================== 1. INTRODUCCIÓN ===========================
\section{Introducción}

El ajedrez fue el problema de referencia de la inteligencia artificial
desde \citet{shannon1950}, y el aprendizaje por refuerzo mostró que una
máquina puede aprender a evaluar posiciones solo jugando: TD-Gammon en el
backgammon \citep{tesauro1995}, KnightCap \citep{baxter2000} y Giraffe
\citep{lai2015} en el ajedrez y, con un cómputo inalcanzable para nosotros,
AlphaZero \citep{silver2018}. Este trabajo se hace dos preguntas: qué se
puede lograr con RL y \emph{una laptop}, y cómo demostrar honestamente que
la mejora se debe al aprendizaje y no a otra cosa.

El caso de aplicación es un robot que juega contra una persona en un tablero
real: ve las jugadas con una cámara, muestra las suyas en una pantalla y
expresa emociones. Esas emociones no son decorativas: la \emph{confianza}
es el valor que la red le da a la posición y la \emph{sorpresa} es la misma
señal de error con la que aprende. El objetivo de fuerza era de unos 1000
puntos de Elo.

Los aportes del trabajo son:
\begin{enumerate}
\item Un sistema de RL completo que entrena en horas en una laptop.
\item Una forma de medir que atribuye la mejora al RL, con intervalos de
  confianza.
\item El diagnóstico de por qué el aprendizaje se estancó y una corrección
  que funcionó (+148 Elo).
\item Un análisis de explicabilidad que muestra qué conocimiento de ajedrez
  agregó el RL.
\item La integración en un robot con cámara, pantalla, reloj de ajedrez y
  emociones que salen del propio RL.
\end{enumerate}

% =========================== 2. MARCO TEÓRICO ===========================
\section{Marco teórico}

\subsection{El ajedrez como problema de RL}

En RL un agente aprende a decidir a partir de recompensas. En el ajedrez,
la única recompensa llega al final: $z=+1$ si gana, $0$ si empata y $-1$ si
pierde. Formalmente es un proceso de decisión de Markov
\citep{sutton2018}: el estado $s$ es la posición y las acciones son las
jugadas legales. Lo que queremos aprender es el \emph{valor} de cada
posición, es decir, qué resultado se espera desde ahí. Lo aproximamos con una
red $\vhat(s)\in[-1,1]$, siempre desde el punto de vista del bando que
mueve, y la leemos como una chance de ganar $p=(\vhat+1)/2$.

\subsection{Aprender de la jugada siguiente: TD($\lambda$)}

Esperar al final de la partida para corregir cada evaluación es lento. La
idea de las diferencias temporales \citep{sutton1988} es corregir cada
predicción con la siguiente: si en la jugada $t$ la red decía $w_t$ y en la
$t+1$ dice $w_{t+1}$, la diferencia $\delta_t=w_{t+1}-w_t$ es una
``sorpresa'' que indica cuánto se equivocó. TD($\lambda$) mezcla todos los
horizontes, desde la jugada siguiente hasta el resultado final:
\begin{equation}
G_t^\lambda = (1-\lambda)\,w_{t+1} + \lambda\,G_{t+1}^\lambda,
\qquad G_T^\lambda = z,
\label{eq:lambda}
\end{equation}
y mueve los pesos hacia ese objetivo:
$\theta \leftarrow \theta + \alpha\,(G_t^\lambda - \vhat(s_t))\,\nabla_\theta \vhat(s_t)$.
Con $\lambda=0$ cada posición aprende solo de la siguiente; con
$\lambda=1$, solo del resultado. Lo importante es que $z$ es un \emph{ancla
externa}: información sobre quién ganó que no sale de la propia red.

\subsection{Mirar hacia adelante: búsqueda}

Una red sola juega mal, porque no ve las consecuencias concretas de una
jugada. Por eso el agente \emph{busca}: prueba sus jugadas, las respuestas
del rival, y se queda con la que mejor resultado garantiza (negamax):
\begin{equation}
v_d(s)=\max_{a}\, -v_{d-1}(s_a), \qquad v_0(s)=\vhat(s),
\end{equation}
donde $d$ es la profundidad y $s_a$ la posición tras la jugada $a$. La poda
alfa-beta \citep{knuth1975} da el mismo resultado sin mirar ramas que no
pueden cambiar la decisión. Para no cortar la búsqueda en medio de un
cambio de piezas (el \emph{efecto horizonte}), en las hojas se siguen las
capturas hasta llegar a una posición tranquila (búsqueda de
\emph{quiescencia}).

\subsection{Dos formas de aprender buscando}

\textbf{TD-Leaf($\lambda$)} \citep{baxter2000} aplica TD($\lambda$) a un
agente que busca: en vez de corregir la posición del tablero, corrige la
posición del final de la variante que eligió la búsqueda, que es la que
realmente evaluó la red.

\textbf{TreeStrap} \citep{veness2009} usa la búsqueda como maestra: cada
posición del árbol aprende el valor que le dio la búsqueda,
\begin{equation}
\mathcal{L}_{\text{TS}}(\theta)=\sum_{s\in\mathcal{T}}\big(v_d(s)-\vhat(s)\big)^2 .
\end{equation}
Aprovecha muchas posiciones por jugada, pero la red aprende de sí misma:
si la red dice siempre 0, la búsqueda también dice 0 y no hay error. Sin un
ancla externa puede caer en ese punto fijo, y lo observamos
(\cref{sec:cero}). Por eso combinamos las dos señales: TreeStrap aporta
cantidad y TD-Leaf aporta el ancla del resultado.

\subsection{Cómo medimos}

Si un jugador saca un puntaje medio $s$ (victoria 1, tablas 0,5) contra
otro, la diferencia de Elo es $D=400\log_{10}\!\big(s/(1-s)\big)$
\citep{elo1978}: +100 Elo equivale a sacar un 64\,\% de los puntos. Para la
predicción del resultado usamos el \emph{AUC}: la probabilidad de que la red
le dé más valor a una posición del bando que terminó ganando que a una del
que terminó perdiendo (0,5 es azar y 1 es perfecto). El AUC no depende de la
escala de la salida: si sube, la red ordena mejor las posiciones, no
solo se volvió más segura.

% =========================== 3. SISTEMA ================================
\section{El sistema de aprendizaje}

\paragraph{Qué ve la red.} La posición entra como 12 tableros de $8\times8$
con unos y ceros, uno por tipo de pieza y color (``¿hay un caballo propio
en esta casilla?''), siempre desde el bando que mueve. La red es una ResNet
chica \citep{he2016} de unos 369.000 parámetros con salida $\tanh$. No ve el
derecho a enrocar, la captura al paso ni las repeticiones (limitación
conocida).

\paragraph{Cómo busca.} A profundidad 1 o 2, con quiescencia en las hojas.
Para que la profundidad 2 entrara en el tiempo disponible diseñamos una
poda alfa-beta compatible con evaluar las posiciones en lote en la GPU:
da el mismo resultado que el árbol completo evaluando unas 250 posiciones
en vez de 1.070 (detalle en el \cref{anx:tecnico}). En las hojas también se
detectan los mates en una jugada.

\paragraph{Cómo aprende.} Cada partida arranca de una posición real de
Lichess \citep{lichess} para tener variedad. El agente juega contra sí
mismo (60\,\%), contra versiones viejas suyas (15\,\%, para no olvidar) y
contra Stockfish \citep{stockfish} (25\,\%), con un currículo que sube o baja
la dificultad según cómo le va \citep{bengio2009}. De cada partida salen
muestras de TreeStrap y de TD-Leaf ($\lambda=0,7$) que se guardan en una
memoria de repetición \citep{lin1992} y se usan para entrenar la red.
Varios procesos juegan en paralelo en la CPU mientras otro entrena en la
GPU; 150.000 partidas a profundidad 1 llevan unas 12 horas en una laptop con
una RTX 4050. Los hiperparámetros están en el \cref{anx:hiper}.

% =========================== 4. DISEÑO EXPERIMENTAL ====================
\section{Diseño experimental}

La pregunta central es de \emph{atribución}: ¿el agente juega mejor porque
cambiaron los pesos aprendiendo de partidas? Para responderla, en cada
comparación todo lo demás queda fijo: misma búsqueda, mismas posiciones,
mismo rival. \Cref{tab:corridas} resume las corridas.

\begin{table}[h]
\centering\small
\begin{tabularx}{\textwidth}{@{}l X r@{}}
\toprule
Corrida & Descripción & Partidas \\
\midrule
\texttt{desde\_cero} & TD-Leaf desde pesos al azar, profundidad 1 & 0 a 20.000 \\
\texttt{treestrap\_cero} & TreeStrap solo, desde pesos al azar & 0 a 20.000 \\
\texttt{final} & Desde una red que conoce el valor del material; objetivo mixto & 0 a 150.000 \\
\texttt{final\_sf} & Continuación de \texttt{final} con Stockfish, profundidad 1 & 150.000 a 300.000 \\
\texttt{mejoras} & Continuación de \texttt{final} con las correcciones de la \cref{sec:meseta}, profundidad 2 & 150.000 a 210.000 \\
\bottomrule
\end{tabularx}
\caption{Corridas de entrenamiento. La red de material, punto de partida de
\texttt{final}, se obtiene con una regresión supervisada del valor de las
piezas, como en KnightCap y Giraffe; todo lo posterior es RL.}
\label{tab:corridas}
\end{table}

Usamos tres instrumentos:
\begin{enumerate}
\item \textbf{Escalera de rivales fijos} (desde un jugador al azar hasta
  Stockfish limitado a 1320 de Elo). Da un Elo absoluto, pero es poco
  precisa: con 100 partidas el puntaje tiene un margen de $\pm0,1$.
\item \textbf{Enfrentamiento directo} entre dos redes con la misma búsqueda:
  500 posiciones de partida reales, cada una jugada dos veces con los colores
  cambiados (1.000 partidas). Detecta diferencias de unos 30 Elo.
\item \textbf{Predicción del resultado} en 47.484 posiciones de partidas de
  Lichess de febrero de 2013, que ninguna red vio (el entrenamiento usó
  enero). Medimos el AUC con intervalos por bootstrap
  \citep{efron1993}.
\end{enumerate}

% =========================== 5. RESULTADOS =============================
\section{Resultados}

\subsection{Desde cero: el RL aprende a evaluar}
\label{sec:cero}

Con pesos al azar la red no sabe ni cuánto vale una pieza. Con TD-Leaf, en
20.000 partidas contra sí misma pasa de sacar 0,05 a 0,85 contra un jugador
que siempre come lo que más gana, y su AUC sube de 0,48 (azar) a 0,73, cerca
de un predictor que solo cuenta material (0,775), algo que nunca se le
enseñó (\cref{fig:cero}). TreeStrap solo, en cambio, colapsa: termina
diciendo 50\,\% en cualquier posición (AUC 0,53). Es el punto fijo de la
sección 2.4 visto en la práctica, y la razón del objetivo mixto.

\begin{figure}[h]
\centering
\includegraphics[width=0.68\textwidth]{desde_cero_auc.jpeg}
\caption{Mejora del AUC respecto de la red al azar (IC 95\,\%). En azul,
TD-Leaf desde cero a 5.000, 10.000 y 20.000 partidas; en gris, el predictor
de material y TreeStrap solo, que colapsó.}
\label{fig:cero}
\end{figure}

\subsection{Desde el material: el RL suma fuerza}

Como arrancar desde cero es lento, la corrida principal parte de la red de
material y deja que el RL cambie la pregunta de ``cuánto material hay'' a
``quién gana''. Tras 150.000 partidas, contra Stockfish 1320 saca 0,385 en
300 partidas, es decir \textbf{unos 1240 de Elo [1200; 1280]}, por encima
del objetivo de 1000 (la red sin RL rondaba los 1020). Con los instrumentos
precisos la atribución es clara: en enfrentamiento directo le gana a la red
sin RL por \textbf{+36 Elo [+26; +47]} y predice mejor quién gana
(\cref{tab:atribucion}).

\subsection{La meseta y su diagnóstico}
\label{sec:meseta}

Seguimos entrenando 150.000 partidas más (\texttt{final\_sf}) y la mejora
fue chica: +31 Elo, con un AUC que quedó igual en 200.000, 250.000 y 300.000
partidas. Revisando el código y los registros encontramos tres causas
principales (el diagnóstico completo está en el \cref{anx:tecnico}):
\begin{itemize}
\item \textbf{Una maestra poco profunda.} A profundidad 1, la búsqueda que
  le enseña a la red ve una sola jugada más que la red: ya no tenía nada
  nuevo para enseñarle.
\item \textbf{Ceguera a las amenazas.} El 73\,\% de las partidas contra
  Stockfish terminaban con el agente recibiendo un mate que no veía.
\item \textbf{Un rival sin escalones.} El nivel más fácil de Stockfish ya
  era demasiado fuerte, así que el currículo nunca avanzó.
\end{itemize}
Las correcciones (corrida \texttt{mejoras}) fueron: entrenar a profundidad
2, detectar mates, agregar niveles intermedios de Stockfish, explorar menos
y ampliar la memoria de la red.

\subsection{Las correcciones rompen la meseta}

\begin{table}[h]
\centering\small
\begin{tabular}{@{}lcc@{}}
\toprule
Enfrentamiento directo (1.000 partidas) & Elo [IC 95\,\%] & AUC de la red nueva \\
\midrule
RL 150k contra sin RL & +36 [+26; +47] & 0,782 (sin RL: 0,767) \\
Profundidad 1: 300k contra 150k & +31 [+20; +43] & 0,790 \\
\texttt{mejoras}: 180k contra 150k & +114 [+100; +127] & 0,799 \\
\texttt{mejoras}: 210k contra 180k & +23 [+9; +36] & 0,801 \\
\textbf{\texttt{mejoras}: 210k contra 150k} & \textbf{+148 [+134; +163]} & \textbf{0,801} \\
\bottomrule
\end{tabular}
\caption{Atribución al RL: en cada fila las dos redes usan la misma
búsqueda, así que la diferencia está en los pesos. Todas las diferencias
son significativas.}
\label{tab:atribucion}
\end{table}

\begin{figure}[t]
\centering
\includegraphics[width=\textwidth]{resumen_atribucion.png}
\caption{Izquierda: diferencia de Elo en enfrentamiento directo (IC 95\,\%).
Derecha: mejora del AUC respecto de la red sin RL. Gris: sin RL; azul: tras
el RL; naranja: continuación a profundidad 1; verde: corrida con las
correcciones.}
\label{fig:atribucion}
\end{figure}

Con las correcciones, 30.000 partidas valen \textbf{+114 Elo}, más que las
150.000 de la continuación a profundidad 1 (\cref{tab:atribucion,fig:atribucion}).
Seguimos hasta 210.000: la fuerza subió unos 25 Elo más hasta las 195.000 y
después se estabilizó (las redes de 195k y 210k le sacan lo mismo a la de
180k), y la predicción ya se había estabilizado en 180k. Por eso cerramos el
entrenamiento en 210.000 partidas, que es la red que usa el robot:
\textbf{+148 Elo} sobre la de 150k.

Revisamos que no hubiera trampas en la medición: el color no influye,
ninguna partida llegó al tope de jugadas y la red con las correcciones
gana más partidas por mate. Y una lección: 100 partidas contra un rival
fijo no detectaban diferencias de 30 a 40 Elo que el enfrentamiento directo
sí detecta.

\subsection{Qué aprendió la red: explicabilidad}
\label{sec:xai}

Las mediciones anteriores muestran \emph{que} el RL mejoró la red. Para ver
\emph{qué} aprendió usamos \emph{oclusión} \citep{zeiler2014}: sacamos una
pieza del tablero y medimos cuánto pierde su dueño según la red,
\[
\text{importancia}(x) = u(s) - u(s\setminus x),
\]
con $u$ el valor para el dueño de la pieza $x$, expresado en peones. La
elegimos en lugar de métodos de gradientes (como Grad-CAM) porque la entrada
es de unos y ceros: un gradiente responde qué pasaría con ``media pieza'',
algo que no existe, mientras que sacar la pieza es una pregunta natural.
Promediamos sobre 3.000 posiciones tranquilas de las partidas de febrero
(sin sacar los reyes).

\begin{table}[h]
\centering\small
\begin{tabular}{@{}lcccc@{}}
\toprule
Valor para la red (en peones) & Sin RL & RL 150k & \texttt{mejoras} 210k & Clásico \\
\midrule
Caballo / alfil & 1,6 / 1,7 & 1,4 / 1,6 & 2,3 / 2,7 & 3 / 3 \\
Torre / dama & 3,0 / 5,3 & 2,8 / 6,4 & 4,6 / 9,0 & 5 / 9 \\
Peón a un paso de coronar & 0,9 & 1,7 & 2,5 & -- \\
\quad el mismo, en peones comunes de cada red & 1,5 & 3,4 & 3,4 & -- \\
\bottomrule
\end{tabular}
\caption{Valor medio de cada pieza para la red, por oclusión, en 3.000
posiciones no vistas. Todos los intervalos al 95\,\% son menores que
$\pm0,1$, salvo los del peón a un paso de coronar (por ejemplo, 0,76 a
0,96 sin RL y 1,5 a 3,4 en 210k).}
\label{tab:xai}
\end{table}

\begin{figure}[t]
\centering
\includegraphics[width=\textwidth]{valores_piezas.png}
\caption{Izquierda: valor de cada pieza para tres redes y los valores
clásicos. Derecha: valor del peón según cuánto avanzó (IC 95\,\%).}
\label{fig:xai_global}
\end{figure}

\Cref{tab:xai,fig:xai_global} muestran dos cosas. Primero, tras el RL los
valores de las piezas se acercan a los clásicos (3, 3, 5 y 9). Una parte de
eso es que la red se volvió más segura de sí misma en general, pero no es
solo eso. Segundo, y más claro: \textbf{el RL aprendió que un peón avanzado
vale mucho más}. Un peón a un paso de coronar vale 1,5 peones comunes para la
red sin RL y 3,4 para la red con RL. La red sin RL partió del material, que
cuenta igual a un peón en cualquier casilla: ese conocimiento solo pudo
venir de jugar.

\begin{figure}[t]
\centering
\includegraphics[width=\textwidth]{mapa_lichess_1.png}
\caption{Oclusión en una posición de Lichess: el número sobre cada pieza es
cuánto pierde su dueño si se la saca; a la derecha, cuánto más la valora la
red tras el RL. Para la red de 210k, el peón negro de d3, a dos pasos de
coronar y apoyado por la dama, vale 12 peones (0,9 para la red sin RL), y
por eso evalúa la posición a favor de las negras. Es un caso extremo,
elegido entre 1.000 posiciones por ser de los que más cambian.}
\label{fig:xai_mapa}
\end{figure}

Los mapas sobre el tablero (\cref{fig:xai_mapa}) explican decisiones
concretas, y son los mismos que el robot muestra con su botón
``¿Por qué?'' (\cref{sec:robot}). La oclusión tiene límites: mide cada
pieza por separado y no sus combinaciones, y al sacar una pieza la posición
puede quedar rara (una dama que de pronto queda atacada), así que los mapas
de posiciones sueltas son ilustrativos y las conclusiones salen de los
promedios.

% =========================== 6. CASO DE APLICACIÓN =====================
\section{Caso de aplicación: el robot}
\label{sec:robot}

El agente juega contra una persona en un tablero de ajedrez real. El
sistema tiene tres partes: \emph{percepción} (ver la jugada de la persona),
\emph{decisión} (elegir la respuesta) y \emph{expresión} (mostrar la jugada y
el estado del agente).

\paragraph{Montaje.} Una webcam montada en un trípode mira el tablero desde
arriba, a aproximadamente un metro de altura, y ve el tablero completo. Una
pantalla conectada a la computadora muestra las jugadas del robot y una
cara que expresa su estado. Como el robot no tiene brazo, la persona mueve
las piezas del robot siguiendo lo que indica la pantalla.
% Si suben la foto del montaje como figuras/montaje.jpg, aparece sola:
\IfFileExists{figuras/montaje.jpg}{%
\begin{figure}[h]
\centering
\includegraphics[width=0.6\textwidth]{montaje.jpg}
\caption{El montaje: webcam en un trípode sobre el tablero y pantalla con
las jugadas y la cara del robot.}
\label{fig:montaje}
\end{figure}}{}

\paragraph{Percepción.} Al empezar se marcan las cuatro esquinas del
tablero en la imagen; con eso la imagen se endereza a una vista cenital y se
divide en 64 casillas. Para detectar una jugada se compara cada casilla
con una foto de referencia y se elige la jugada \emph{legal} que mejor
explica las casillas que cambiaron: conocer las reglas filtra el ruido. Con
la cámara real aparecían cambios falsos (la cámara reajusta sola el brillo,
algunas luces parpadean y el trípode se puede mover un poco). Lo resolvimos
normalizando el brillo, tomando la mediana de 15 imágenes quietas y
comparando de forma tolerante a corrimientos de uno o dos píxeles. En una
simulación con corrimientos, el detector original acertaba 12 de 23
jugadas y el nuevo las 23, sin detectar nunca una jugada equivocada. Con
el montaje real jugamos varias partidas completas y funcionó bien.

\paragraph{Reloj de ajedrez.} La persona avisa que terminó su jugada con la
barra espaciadora, que funciona como el botón de un reloj de ajedrez real:
cada jugador tiene un tiempo (por ejemplo, 5 minutos más 3 segundos por
jugada) y corre el reloj de quien tiene que mover. Si la cámara confirma
una jugada legal, el reloj pasa al rival; si no, sigue corriendo. El reloj
del robot corre mientras piensa y hasta que su jugada se hace en el
tablero. Quien se queda sin tiempo pierde, salvo que al rival no le
alcance el material para dar mate (tablas), como en el reglamento.

\paragraph{Decisión.} El robot usa la red de 210.000 partidas con búsqueda
a profundidad 2 y detección de mates, la misma con la que entrenó. Su
jugada queda pendiente hasta que la cámara confirma que se hizo en el
tablero.

\paragraph{Expresión.} La pantalla muestra la jugada del robot con una
flecha sobre un tablero digital y en palabras (``Caballo come en e5,
jaque''), y una cara que sale de dos señales del propio RL. La
\emph{confianza} es la chance de ganar $p=(v+1)/2$, y la \emph{sorpresa} es
el error TD entre dos jugadas seguidas, $\delta=v_t-v_{t-1}$. Si
$|\delta|>0,16$ la cara se muestra \emph{sorprendida}; si no, según $p$:
\emph{ganando} ($p\ge0,75$), \emph{confiada} ($p\ge0,5$),
\emph{preocupada} ($p\ge0,3$) o \emph{perdiendo}. Así, lo que el público ve
en la cara del robot es literalmente lo que el agente usa para aprender.

\paragraph{¿Por qué?} Un botón muestra sobre el tablero digital el mapa de
oclusión de la posición actual (\cref{sec:xai}): cuánto vale para el robot
cada pieza en ese momento. Así la persona puede ver, por ejemplo, que el
robot está preocupado porque le da mucho valor a un peón rival que está por
coronar.

% =========================== 7. DISCUSIÓN ==============================
\section{Discusión y limitaciones}

\begin{itemize}
\item \textbf{Las correcciones se aplicaron juntas.} Por el diagnóstico,
  atribuimos la mayor parte a la profundidad 2, pero no corrimos el
  experimento que separaría el aporte de cada corrección.
\item \textbf{Una sola semilla.} Por costo de cómputo cada corrida se
  entrenó una vez: los intervalos cubren el ruido de la medición, no el del
  entrenamiento.
\item \textbf{Elo en escala Stockfish.} Los 1240 son de la red de 150k contra
  Stockfish y no equivalen a un rating humano; el Elo absoluto de la red de
  210k no se midió (le gana a la de 150k por +148).
\item \textbf{Entrada incompleta.} La red no ve el enroque, la captura al
  paso ni las repeticiones.
\item \textbf{Nueva meseta.} La corrida con correcciones se estabilizó hacia
  las 200.000 partidas: para seguir mejorando harían falta cambios de fondo
  (más información de entrada o una red más grande), no solo más partidas.
\item \textbf{Visión y explicabilidad.} Las pruebas sistemáticas de la
  visión son simuladas; las reales fueron partidas completas sin un conteo
  formal de aciertos. La oclusión mide piezas por separado y los mapas
  sueltos son ilustrativos.
\end{itemize}

% =========================== 8. CONCLUSIONES ===========================
\section{Conclusiones y trabajo futuro}

Mostramos que el RL basado en valor, combinado con búsqueda, entrena un
jugador de ajedrez razonable en una laptop, y que la mejora se puede
atribuir al aprendizaje con mediciones cuidadosas. El recorrido tuvo cinco
pasos: (1) desde cero, el RL aprende a evaluar casi como el material sin
que nadie se lo enseñe; (2) desde el material, suma +36 Elo; (3) seguir
igual se estanca; (4) diagnosticamos por qué, lo corregimos y la red llega
a +148 Elo antes de estabilizarse otra vez; (5) la explicabilidad muestra
que esa mejora tiene contenido de ajedrez real, como el valor de los peones
avanzados. La lección más general es sobre cómo medir: con 100 partidas
contra un rival fijo no se veía nada; con enfrentamientos directos y
predicción fuera de muestra, sí.

Como trabajo futuro: medir el Elo absoluto de la mejor red, separar el
aporte de cada corrección, darle a la red la información que le falta o una
red más grande para superar la nueva meseta, buscar conceptos más complejos
dentro de la red (rey expuesto, estructura de peones) y sumarle un brazo al
robot.

% =========================== BIBLIOGRAFÍA ==============================
\clearpage
\begin{thebibliography}{99}
\setlength{\itemsep}{2pt}\small

\bibitem[Baxter et~al.(2000)]{baxter2000}
Baxter, J., Tridgell, A. y Weaver, L. (2000). Learning to play chess using
temporal differences. \textit{Machine Learning}, 40(3), 243--263.

\bibitem[Bengio et~al.(2009)]{bengio2009}
Bengio, Y., Louradour, J., Collobert, R. y Weston, J. (2009). Curriculum
learning. \textit{Proceedings of ICML}.

\bibitem[Brier(1950)]{brier1950}
Brier, G. W. (1950). Verification of forecasts expressed in terms of
probability. \textit{Monthly Weather Review}, 78(1), 1--3.

\bibitem[Efron y Tibshirani(1993)]{efron1993}
Efron, B. y Tibshirani, R. J. (1993). \textit{An Introduction to the
Bootstrap}. Chapman \& Hall.

\bibitem[Elo(1978)]{elo1978}
Elo, A. E. (1978). \textit{The Rating of Chessplayers, Past and Present}.
Arco.

\bibitem[He et~al.(2016)]{he2016}
He, K., Zhang, X., Ren, S. y Sun, J. (2016). Deep residual learning for
image recognition. \textit{Proceedings of CVPR}.

\bibitem[Huber(1964)]{huber1964}
Huber, P. J. (1964). Robust estimation of a location parameter.
\textit{Annals of Mathematical Statistics}, 35(1), 73--101.

\bibitem[Kingma y Ba(2015)]{kingma2015}
Kingma, D. P. y Ba, J. (2015). Adam: a method for stochastic optimization.
\textit{Proceedings of ICLR}.

\bibitem[Knuth y Moore(1975)]{knuth1975}
Knuth, D. E. y Moore, R. W. (1975). An analysis of alpha-beta pruning.
\textit{Artificial Intelligence}, 6(4), 293--326.

\bibitem[Lai(2015)]{lai2015}
Lai, M. (2015). \textit{Giraffe: using deep reinforcement learning to play
chess}. Tesis de maestría, Imperial College London.

\bibitem[Lichess(2013)]{lichess}
Lichess (2013). \textit{Lichess open database}, partidas de enero y febrero
de 2013. \url{https://database.lichess.org}.

\bibitem[Lin(1992)]{lin1992}
Lin, L.-J. (1992). Self-improving reactive agents based on reinforcement
learning, planning and teaching. \textit{Machine Learning}, 8, 293--321.

\bibitem[Shannon(1950)]{shannon1950}
Shannon, C. E. (1950). Programming a computer for playing chess.
\textit{Philosophical Magazine}, 41(314), 256--275.

\bibitem[Silver et~al.(2018)]{silver2018}
Silver, D. et~al. (2018). A general reinforcement learning algorithm that
masters chess, shogi, and Go through self-play. \textit{Science}, 362(6419),
1140--1144.

\bibitem[Stockfish(s.f.)]{stockfish}
Stockfish. Motor de ajedrez de código abierto.
\url{https://stockfishchess.org}.

\bibitem[Sutton(1988)]{sutton1988}
Sutton, R. S. (1988). Learning to predict by the methods of temporal
differences. \textit{Machine Learning}, 3(1), 9--44.

\bibitem[Sutton y Barto(2018)]{sutton2018}
Sutton, R. S. y Barto, A. G. (2018). \textit{Reinforcement Learning: An
Introduction}, 2.ª ed. MIT Press.

\bibitem[Tesauro(1995)]{tesauro1995}
Tesauro, G. (1995). Temporal difference learning and TD-Gammon.
\textit{Communications of the ACM}, 38(3), 58--68.

\bibitem[Veness et~al.(2009)]{veness2009}
Veness, J., Silver, D., Uther, W. y Blair, A. (2009). Bootstrapping from
game tree search. \textit{Advances in NeurIPS}, 22.

\bibitem[Wilson(1927)]{wilson1927}
Wilson, E. B. (1927). Probable inference, the law of succession, and
statistical inference. \textit{Journal of the American Statistical
Association}, 22(158), 209--212.

\bibitem[Zeiler y Fergus(2014)]{zeiler2014}
Zeiler, M. D. y Fergus, R. (2014). Visualizing and understanding
convolutional networks. \textit{Proceedings of ECCV}, 818--833.

\end{thebibliography}

% =========================== ANEXOS ====================================
\appendix
\clearpage
\section{Hiperparámetros}
\label{anx:hiper}

\begin{table}[h]
\centering\footnotesize
\begin{tabularx}{\textwidth}{@{}>{\raggedright\arraybackslash}p{0.3\textwidth}>{\raggedright\arraybackslash}X>{\raggedright\arraybackslash}X@{}}
\toprule
Parámetro & \texttt{final} / \texttt{final\_sf} & \texttt{mejoras} \\
\midrule
Profundidad de búsqueda / quiescencia & 1 / 4 & 2 con poda / 4, mates en 1 \\
Algoritmo & mixto (TreeStrap + TD-Leaf) & igual \\
$\lambda$ / repeticiones de cada muestra TD & 0,7 / 4 & igual \\
Muestras de TreeStrap por jugada & 8 (raíz + 7 al azar) & igual \\
Temperatura de exploración & 0,05 en todas las jugadas & 0,05 en las 10 primeras; 0 contra Stockfish \\
Largo del autojuego & tramos de 30 jugadas & 20\,\% hasta el final (máx.\ 200) \\
Rivales & autojuego 60\,\%, viejos 15\,\%, Stockfish 25\,\%$^{*}$ & igual, con 6 escalones al azar \\
Memoria / lote / usos por muestra & 400.000 / 512 / 2,0 & 2.000.000 / 512 / 2,0 \\
Optimizador \citep{kingma2015} & Adam, $5\cdot10^{-4}$, decaimiento $10^{-5}$ & tasa lineal de $5\cdot10^{-4}$ a $10^{-4}$ \\
Pérdida \citep{huber1964} & Huber ($\beta=0,1$) & igual \\
Procesos que juegan & 7 (RTX 4050) / 6 (Apple M1) & 7 (RTX 4050) \\
\bottomrule
\end{tabularx}
\caption{Hiperparámetros principales. $^{*}$\texttt{final} no usó Stockfish
(autojuego y versiones viejas).}
\end{table}

\section{Detalles técnicos}
\label{anx:tecnico}

\paragraph{Búsqueda en lote.} Todas las hojas se evalúan juntas, de modo
que la red corre una vez por jugada. La quiescencia de las hojas usa solo
material (P 100, C 320, A 330, T 500, D 900 centipeones) con poda delta, y
la posición tranquila resultante es la que va a la red. La poda alfa-beta a
profundidad 2 tiene tres pasos: se explora completa la mejor jugada según
profundidad 1; de las demás se evalúan en un lote las 4 primeras
respuestas; y las que no quedan refutadas se siguen por tandas. En 150
posiciones reales da el mismo valor de la raíz que el árbol completo
(diferencia máxima $2\cdot10^{-7}$).

\paragraph{Experiencia y objetivo mixto.} Cada partida arranca de una
posición de Lichess de enero de 2013 más una jugada al azar. El autojuego se
juega en tramos de 30 jugadas; si no termina, se adjudica por material: la
ventaja $c$ en centipeones se lleva a $z=\tanh(c\ln10/800)$. En cada jugada
se guardan hasta 8 muestras de TreeStrap y, al terminar, una de TD-Leaf por
jugada, repetida 4 veces para equilibrar su peso. La pérdida es Huber,
robusta a los objetivos atípicos de TD.

\paragraph{Diagnóstico completo de la meseta.} Además de las tres causas de
la \cref{sec:meseta}: el currículo medía al agente con exploración, el
88\,\% del autojuego se adjudicaba por material (algo que la red ya sabía)
y la memoria de 400.000 muestras equivalía a unos 4 minutos de juego. Las
correcciones completas: profundidad 2 con poda, mates en 1, 6 escalones de
Stockfish que juegan al azar con probabilidad 0,8 a 0,1, currículo medido
sin exploración, exploración solo en las 10 primeras jugadas, 20\,\% del
autojuego hasta el final, memoria de 2 millones de muestras y una tasa de
aprendizaje que baja de $5\cdot10^{-4}$ a $10^{-4}$.

\paragraph{Mediciones.} El enfrentamiento directo arranca de posiciones
reales entre las jugadas 5 y 15 con material parejo; como las redes son
deterministas, eso da variedad, y el intervalo se calcula por par de
partidas. Control: una red contra sí misma da 0,500. En la predicción, los
cambios de piezas pendientes se resuelven con la quiescencia y el bootstrap
es pareado por partida: lo que importa es el intervalo de la diferencia
entre redes. También medimos el error de Brier \citep{brier1950}, que
mejora en forma significativa ($-0,007$ para 180k y 210k), y los puntajes
de la escalera se reportan con el intervalo de Wilson \citep{wilson1927}.
Desde las 190.000 partidas la red se vuelve más segura de sus valores sin
ordenarlos mejor: el AUC no cambia y el LogLoss empeora levemente.

\paragraph{Oclusión.} El valor $v$ se pasa a peones con la escala del
preentrenamiento de la red de material:
$c=800/\ln 10\cdot\operatorname{atanh}(v)$ centipeones. Las posiciones son
tranquilas (sin jaque y sin capturas pendientes) y se descartan las que
quedan ilegales al sacar una pieza. Los intervalos son por bootstrap de
posiciones. La red de material ordena bien las piezas pero con valores
encogidos hacia 0 (efecto del error cuadrático sobre posiciones casi
parejas); tras el RL los cocientes entre piezas también cambian, así que el
cambio no es solo de escala.

\section{Calibración de las redes}
\label{anx:calib}

\begin{figure}[h]
\centering
\includegraphics[width=0.6\textwidth]{valor_calibracion.jpeg}
\caption{Proporción real de puntos frente a la chance de ganar que da cada
red, en las posiciones de febrero de 2013. Todas las redes están cerca de la
diagonal; el predictor de material puro es el más alejado en las posiciones
desfavorables.}
\end{figure}

\section{Reproducción}

El código está en el repositorio del proyecto (Python con PyTorch y
\texttt{python-chess}, sin notebooks). Comandos principales:

{\small
\begin{verbatim}
# Red de material y corrida principal
python preentrenar_material.py --out checkpoints/material.pt --fens data/fens.txt
python entrenar.py --out runs/final --init checkpoints/material.pt --algo mixto \
    --start lichess --fens data/fens.txt --workers 7 --games 150000

# Corrida con las correcciones (hasta 210.000 partidas)
python entrenar.py --out runs/mejoras --init runs/final/ck_150000.pt --resume \
    --algo mixto --start lichess --fens data/fens.txt --stockfish RUTA \
    --workers 7 --games 60000 --depth 2 --poda --mates --sf-azar \
    --temp-stockfish 0 --temp-jugadas 10 --p-completas 0.2 \
    --buffer 2000000 --lr-final 1e-4

# Mediciones de atribución (todas juntas: bash lanzar_mediciones.sh)
python medir_valor.py --pgn data/lichess_db_standard_rated_2013-02.pgn.zst \
    --ck sin_rl=runs/final/ck_000000.pt rl=runs/final/ck_150000.pt
python duelo.py --a runs/mejoras/ck_210000.pt --b runs/final/ck_150000.pt \
    --pares 500 --poda

# Explicabilidad por oclusión
python explicar.py --ck sin_rl=runs/final/ck_000000.pt \
    rl_150k=runs/final/ck_150000.pt mej_210k=runs/mejoras/ck_210000.pt

# Jugar contra el robot (webcam, o --sin-camara para el simulador)
python chesscam/chesscam.py --ck runs/mejoras/ck_210000.pt --reloj 5+3
\end{verbatim}
}

\end{document}
